\section{Implementation and Software Framework}

This section details the software architecture and the algorithmic implementation of the vacuum grasping pipeline. The framework is implemented in Python and follows a modular design to decouple geometric processing from hardware-specific parameters.

\subsection{Software Architecture and Configuration Management}

The grasping pipeline is implemented in Python, utilizing an object-oriented programming (OOP) architecture. This approach separates the geometric processing, collision evaluation, and hardware definition into distinct software modules. To manage data structures and algorithm states, the system relies on Python \texttt{dataclasses} combined with static type hinting. This design enforces strict structuring of input parameters and minimizes type-related runtime errors.


To ensure adaptability across different industrial assembly tasks and object geometries, the framework avoids hard-coded numerical values in the source code. Pipeline parameterization is handled through external YAML configuration files. Operational variables, such as collision margins, voxel grid resolutions, and Grasp Stability Score (GSS) weights, are imported dynamically during system initialization.

This parametric modularity extends to the geometric configuration of the end-effectors. The physical properties of the vacuum grippers are abstracted within dedicated hardware modules. Parameters including the number of suction cups, dynamic sealing radii, and relative distances between pad centers are defined externally. As a result, the core algorithmic class (\texttt{VacuumGraspSampler}) remains hardware-agnostic. Transitioning between evaluations using a single-cup gripper and an orthogonal four-cup configuration requires only the substitution of the respective configuration file, leaving the raycasting target discovery and evaluation logic unmodified.

\subsection{Core Libraries and Computational Ecosystem}

The computational efficiency of the framework depends on a specialized ecosystem of open-source Python libraries. Each library is selected to handle specific domains of the grasping pipeline, minimizing execution time and memory overhead.

\textbf{Open3D:} Open3D serves as the primary engine for 3D spatial data processing. While the pipeline is scripted in Python, Open3D operates on an optimized C++ backend, allowing the framework to process dense point clouds efficiently. Key functionalities include voxel downsampling, KD-Tree spatial indexing for nearest-neighbor queries, and tensor processing. Specifically, the surface discovery phase utilizes Open3D's tensor-based \texttt{RaycastingScene}. This module enables the parallelized computation of thousands of inward-facing approach vectors simultaneously, bypassing the computational bottlenecks of standard Python operations.

\textbf{NumPy:} Linear algebra and mathematical evaluations within the Grasp Stability Score (GSS) loop are executed using NumPy. This library computes homogeneous transformations, coordinate frame rotations, and vector operations such as dot and cross products. By vectorizing the mathematical logic required for the torque and verticality evaluations, NumPy eliminates the need for iterative Python loops, significantly reducing the computational time required to process and rank candidate poses.

\textbf{Trimesh:} Geometric collision avoidance is managed by the \texttt{trimesh} library. While Open3D handles the environmental point cloud data, \texttt{trimesh} is strictly dedicated to loading, inflating, and querying the physical 3D CAD meshes of the robotic end-effectors. It provides the boolean intersection algorithms required to verify spatial clearances between the inflated gripper volume and the target object, ensuring collision-free approach trajectories prior to hardware execution.


\subsection{Point Cloud Ingestion and Preprocessing}

The first stage of the computational pipeline involves standardizing the input 3D data. Raw point clouds generated from CAD assemblies or captured by depth sensors are frequently encoded in millimeters. To ensure dimensional compatibility with standard robotic simulation environments and kinematic solvers (e.g., ROS and MoveIt), the software applies a uniform scaling factor of $0.001$ relative to the global origin $(0, 0, 0)$ when the scale of the object is in milimeters. This operation explicitly converts all spatial coordinates to meters before any geometric evaluation occurs.

Following spatial normalization, the point cloud undergoes a voxel downsampling process.  This algorithm divides the 3D space into a regular grid of volumetric elements (voxels) and replaces all points within each voxel with their geometric centroid. This operation serves two critical functions. First, it reduces the memory footprint and the total number of spatial queries required during the raycasting phase. Second, it normalizes the spatial density of the point cloud, compensating for areas of high concentration typical of high-resolution 3D scanners. A uniform global point density is essential to ensure that the density-dependent algorithms—specifically the multi-pad seal verification—perform consistently regardless of the initial sensor resolution.

Finally, a radius-based outlier removal filter is applied to the downsampled cloud. This filter removes isolated spatial noise, which is commonly introduced by depth sensor inaccuracies or reflective object surfaces. The algorithm evaluates the local neighborhood of each point and discards those that fail to meet a minimum threshold of neighbors within a specified radius. This filtering process provides a clean structural representation of the target object, preserving the sharp edges and flat planar surfaces strictly required for successful vacuum grasping.


\subsection{Raycasting Engine and Hardware Symmetry}

The target discovery phase avoids heuristic random sampling by employing a deterministic inward raycasting strategy.  This mechanism is implemented using Open3D's tensor-based \texttt{RaycastingScene} module. Approach vectors are formulated as tensor arrays and cast in batches toward the target object. By utilizing tensor data structures, the software achieves highly parallelized computation of surface intersections, avoiding the computational bottlenecks associated with iterative loops in standard Python.

Because the mathematically computed ray intersections do not inherently correspond to the discrete spatial coordinates of the input point cloud, a spatial mapping step is required. The software utilizes a \texttt{KDTree} nearest-neighbor search to snap each valid ray intersection coordinate to the closest exact physical point on the object's surface. This ensures that the generated grasp candidates are grounded on actual sensor data rather than interpolated empty space.

After establishing a base approach vector aligned with the local surface normal, the algorithm generates a set of candidate orientations by rotating the gripper around this local Z-axis. To minimize computational overhead, this rotational search space is dynamically constrained based on the geometric symmetry of the designated hardware configuration. For instance, a single circular suction cup possesses continuous radial symmetry, allowing the algorithm to bypass Z-axis rotation generation entirely ($0^\circ$). Conversely, a bilateral two-pad gripper is evaluated over a bounded $180^\circ$ range, while an orthogonal four-pad configuration is constrained to a $90^\circ$ search space. This hardware-aware optimization reduces the total number of evaluated candidates, lowering processing time without compromising the discovery of valid grasp poses.


\subsection{Collision Detection Engine}

Before computing the Grasp Stability Score (GSS), the framework evaluates physical clearances to guarantee collision-free approach trajectories. This spatial verification is implemented using the \texttt{trimesh} library. The physical CAD mesh of the vacuum gripper is loaded and dynamically inflated by a predefined collision margin.  This volumetric expansion acts as a safety buffer to absorb point cloud noise and mechanical calibration inaccuracies. A boolean intersection query is then executed between this inflated gripper mesh—positioned at the candidate approach pose—and the target object's point cloud. Candidate poses that trigger an intersection are immediately rejected.



\subsection{Vectorized Grasp Stability Evaluation (GSS)}

After the geometric and collision constraints are satisfied, the remaining candidate poses are evaluated using the Grasp Stability Score (GSS). To ensure high-performance execution within the pipeline, the mathematical scoring functions are heavily vectorized using NumPy, reducing the evaluation time per candidate to fractions of a millisecond.

The first two components of the GSS are the Verticality and Torque evaluations. The Verticality score assesses the angular deviation of the approach vector relative to the global gravity vector (the $Z$-axis). This is computed efficiently across all candidates simultaneously using vectorized dot products. Similarly, the Torque evaluation minimizes the induced mechanical moment by calculating the planar distance in the $XY$ coordinate space from each candidate's Target Center Point (TCP) to the object's estimated Center of Mass (CoM).

The third and most critical component is the Multi-pad Seal Verification, which algorithmically simulates the physical deformation limits of the suction cup bellows. For each candidate pose, the conceptual centers of the gripper pads are projected onto the local point cloud surface.  A KD-Tree spatial index performs a rapid radial search to extract all neighboring surface points within the physical radius of each pad.

To guarantee a hermetic seal, it is insufficient to simply count the total number of points within this radius, as the surface might be discontinuous. Instead, the extracted points are partitioned into 8 distinct angular slices of $45^\circ$ each, distributed radially around the local $Z$-axis of the pad. The algorithm then cross-references the point count in each slice against the expected dynamic point density ($\rho$) calculated during the preprocessing phase. If any of the 8 slices fails to meet the minimum point density threshold, the candidate pose is immediately rejected. This rigorous topological verification prevents vacuum leakage over holes, edges, or sparse surface geometries.
